\documentclass[10pt,a4paper,sans]{moderncv}
\moderncvstyle{classic}
\moderncvcolor{green}

\usepackage[scale=0.85]{geometry}
\setlength{\footskip}{136.00005pt}

\ifxetexorluatex
\usepackage{fontspec}
\usepackage{unicode-math}
\defaultfontfeatures{Ligatures=TeX}
\setmainfont{Liberation Serif}
\setsansfont{Liberation Sans}
\setmonofont{Liberation Mono}
\setmathfont{Latin Modern Math}
\else
\usepackage[T1]{fontenc}
\usepackage{lmodern}
\fi

% document language
\usepackage[english, russian]{babel}

% personal data
\name{Максимов}{Матвей Алексеевич}
\title{Системный администратор GNU/Linux}
\born{25 апреля 2005}
\address{Ворошилова, 27--22}{426053, Ижевск}{Россия}
\phone[mobile]{+7~982~119-11-86}
\email{maksimovmatvey2005@gmail.com}
% \email{maximov-matvi@yandex.ru}

% Social icons
\social[github]{MMatveyA}
\social[gitlab]{MMatveyA}
\social[stackoverflow]{23356580/Матвей-Максимов}
\social[telegram]{mmatvey1123}
\social[whatsapp]{89821191186}
\social[matrix]{@mmatvey:matrix.org}

\extrainfo{Это резюме полностью создано в терминале \\
с помощью \LaTeX{} и NeoVim}
% \photo[98pt][0.4pt]{picture}

\begin{document}
\makecvtitle

\section{Образование}

\cventry{2023--2029}{Радиотехнические системы и комплексы}{ИжГТУ
	им.~М.~Т.~Калашникова}{Ижевск}{\textit{Специалитет}}{}

\section{Опыт работы}

\subsection{Профессиональный}

\cventry{2022--2024}{Системный администратор}{Фриланс}{Россия}{}{Эпизодическая
	помощь с настройкой серверов на базе GNU/Linux, установка и настройка простых
	серверных приложений на базе серверов Apache, Nginx; проектирование баз данных
	для MySQL и PostgreSQL}
\subsection{Прочее}
\cventry{2024}{Проводник пассажирского вагона}{АО
	``ФПК''}{Ижевск}{}{Работа проводником в летний период. Увольнение
	в связи с началом нового учебного года}

\section{Основные навыки}

\cvskillplainlegend*[0.2em][Базовые знания][Расширенные знания][Опыт
	использования в проектах][Обширный опыт использования][Эксперт]{}
\cvskillhead[0.1em][Знания][Навык][Опыт][Комментарий]

\cvskillentry*{ОС:}{4}{GNU/Linux}{4}{Использование разнообразных и
	специфичных дистрибутивов}
\cvskillentry*{Базы данных:}{3}{PostgreSQL}{2}{Установка и настройка базы
	данных для обеспечения доступа по локальной сети и сохранения безопасности
	подключения; установка и настройка \href{https://www.pgadmin.org/}{pgAdmin}}
\cvskillentry{}{1}{MySQL}{1}{Установка и базовая настройка; подключение
	\href{https://www.phpmyadmin.net/}{phpMyAdmin}}
\cvskillentry{}{1}{MongoDB}{1}{Установка и базовая настройка}
\cvskillentry*{Языки программирования:}{4}{Bash}{4}{Написание скриптов для
	автоматизации рутинных процессов, активное использование различных
	возможностей языка}
\cvskillentry{}{3}{C++}{5}{Проектирование и написание программ различного
	профиля: серверные программы, обычные пользовательские. Знание различных
	систем автоматизации сборки программ и библиотек на C++}
\cvskillentry{}{3}{Python}{4}{Проектирование и написание различных программ,
	библиотек}
\cvskillentry{}{2}{SQL}{2}{Проектирование баз данных,
	включающих в себя множество связей с помощью ключей (\texttt{foreign key});
	написание запросов к этим базам данных}

\section{Иностранные языки}

\cvitemwithcomment{Английский}{Технический}{Чтение технической
	литературы, базовый разговорный}
\cvitemwithcomment{Японский}{Слабый}{В процессе изучения}

\section{Прочее}

\cvitem{Командная строка}{Уверенное использование различных командных оболочек,
	способность решить большинство возникающих проблем через командную строку.}
\cvitem{GNU/Linux}{Помимо использования простых и не сложных в настройке
	дистрибутивов GNU/Linux (таких, как Debian, Fedora, ArchLinux), есть опыт в
	использовании более специфичных дистрибутивов: Gentoo, Artix, Devuan
	(дистрибутивы, свободные от SystemD; Gentoo также позволяет гибко настраивать
	процесс компиляции пакетов через USE--флаги), Void Linux (дистрибутив,
	использующий \href{https://www.musl-libc.org}{musl-libc} вместо
	\href{https://www.gnu.org/software/libc/}{glibc}), NixOS (дистрибутив,
	нацеленный на воспроизводимость установки и полностью
	настраиваемый через единый
	конфигурационный файл). Сейчас в процессе изучения возможностей использования
	дистрибутивов, основанных на GNU/Linux, на процессорах, использующих
	архитектуру ARMv8.}

\end{document}
